\chapter{La surface de volatilité et le trading de volatilité}\label{ch:iv2:vol-surface-and-vol-trading}

\begin{ivtrains}
Lire une surface, la faire bouger correctement quand le sous-jacent bouge, et trader la volatilité avec la variance, la dispersion et les événements. Sauf indication contraire de l'énoncé, les taux sont nuls, et 252 jours de cotation font une année. Tous les niveaux ici sont inventés pour les calculs.
\end{ivtrains}

\section{Lecture d'une surface}

\begin{question}[$\star$ \iqroles{QT, Trader, QR, Risque} \iqfirm{toutes} \iqlevel{junior}]\label{iq:iv2:vol-surface-and-vol-trading:1}
Les cinq rendements quotidiens d'une action sont $+1~\%$, $-1~\%$, $+2~\%$, $-2~\%$ et $0~\%$. Quelle est sa volatilité réalisée annualisée, en prenant le rendement moyen égal à zéro~? Pourquoi les desks prennent-ils en général la moyenne égale à zéro pour des données quotidiennes~?
\end{question}

\begin{question}[$\star$ \iqroles{QT, Trader, QR, Risque} \iqfirm{toutes} \iqlevel{junior}]\label{iq:iv2:vol-surface-and-vol-trading:2}
Une action présente une volatilité implicite de 24~\%. Quel mouvement attend-on à peu près sur une journée, selon la règle du 16~? Pourquoi 16~? Quelle serait la même valeur si l'on compte chaque jour calendaire~?
\end{question}

\begin{question}[$\star$ \iqroles{QT, Trader, QR} \iqfirm{toutes} \iqlevel{junior}]\label{iq:iv2:vol-surface-and-vol-trading:3}
Vous êtes acheteur d'une option à la monnaie, couverte en delta, sur une action à 100~\$ avec une volatilité implicite de 25~\%, et vous la conservez un jour de cotation. Quelle amplitude de mouvement faut-il pour atteindre le point mort~? À quelle fréquence l'action dépassera-t-elle ce seuil, si le modèle est valide~?
\end{question}

\begin{question}[$\star$ \iqroles{QT, Trader, QR, Risque} \iqfirm{toutes} \iqlevel{junior}]\label{iq:iv2:vol-surface-and-vol-trading:4}
La volatilité implicite à la monnaie est de 25~\% pour trois mois et de 22~\% pour six mois. Quelle est la volatilité forward entre trois et six mois~? Pourquoi fait-on la moyenne des variances et non des volatilités~?
\end{question}

\begin{question}[$\star$ \iqroles{QT, Trader, QR} \iqfirm{teneur de marché, hedge fund, banque} \iqlevel{junior}]\label{iq:iv2:vol-surface-and-vol-trading:5}
Une action à 100~\$ présente une volatilité implicite à la monnaie de 30~\% à trois mois et de 20~\% à six mois. Qu'est-ce qui ne va pas~? Que tradez-vous, et combien cela vous rapporte-t-il d'emblée~?
\end{question}

\begin{question}[$\star$ \iqroles{QT, Trader, QR, Analyste} \iqfirm{toutes} \iqlevel{junior}]\label{iq:iv2:vol-surface-and-vol-trading:6}
Sur un indice actions, un put à un an situé 10~\% sous le niveau de l'indice a une volatilité implicite de 26~\%, l'option à la monnaie a 20~\%, et le call situé 10~\% au-dessus a 17~\%. Pourquoi la surface a-t-elle cette forme alors que le modèle de Black--Scholes prévoit qu'elle soit plate~? Donnez vos trois meilleures raisons, et dites laquelle vous présenteriez en premier.
\end{question}

\begin{question}[$\star$ \iqroles{QT, Trader, QR} \iqfirm{teneur de marché, hedge fund, banque} \iqlevel{junior}]\label{iq:iv2:vol-surface-and-vol-trading:7}
Un swap de variance a un notionnel vega de 100\,000~\$ et un strike de 20. Quel est le notionnel de variance~? Combien gagnez-vous si vous êtes acheteur et que la volatilité réalisée vaut 25~? Et si elle vaut 15~? Pourquoi le gain est-il plus important que la perte pour un même écart de cinq points~?
\end{question}

\begin{question}[$\star$ \iqroles{QR, QT, Risque} \iqfirm{toutes} \iqlevel{junior}]\label{iq:iv2:vol-surface-and-vol-trading:8}
Vous mesurez une volatilité réalisée de 20~\% à partir de 20 rendements quotidiens d'une action. Quelle est à peu près l'erreur type de ce nombre si les rendements sont normaux~? Combien de jours faut-il pour la ramener à 1 point~?
\end{question}

\begin{question}[$\star$ \iqroles{QT, Trader, QR} \iqfirm{hedge fund, teneur de marché, banque} \iqlevel{junior}]\label{iq:iv2:vol-surface-and-vol-trading:9}
Un indice dont de nombreux membres sont pondérés à parts égales présente une volatilité implicite de 18~\%, et chaque membre a une volatilité implicite de 30~\%. Quelle est la corrélation implicite de l'indice~? Que vous indique une hausse de ce chiffre~?
\end{question}

\section{Sourires de volatilité et leurs déplacements}

\begin{question}[$\star\star$ \iqroles{QR, QT} \iqfirm{hedge fund, teneur de marché, banque} \iqlevel{junior, confirmé}]\label{iq:iv2:vol-surface-and-vol-trading:10}
Vous pensez que la volatilité d'une action à 100~\$ pour l'année suivante vaut soit 10~\%, soit 30~\%, avec des probabilités égales, et vous la connaîtrez tout de suite. Évaluez les calls comme la moyenne des deux prix de Black--Scholes. Quelle volatilité implicite cela donne-t-il pour des strikes de 100, 80 et 130~? Pourquoi la volatilité à la monnaie est-elle proche de la moyenne des deux, et non de leur moyenne quadratique~?
\end{question}

\begin{question}[$\star\star$ \iqroles{QR, QT, QD} \iqfirm{toutes} \iqlevel{junior, confirmé}]\label{iq:iv2:vol-surface-and-vol-trading:11}
Pour une action à 100~\$ avec une échéance d'un an et un taux nul, vous observez des calls à 95, 100 et 105 valant 10.52, 7.97 et 5.91 lorsqu'ils sont évalués à une volatilité de 20~\%. Donnez la densité de probabilité que le marché implique en 100 à partir de ces trois prix. Quelle est la véritable densité lognormale en ce point~? Que signifierait une réponse négative~?
\end{question}

\begin{question}[$\star\star$ \iqroles{QT, Trader, QR} \iqfirm{teneur de marché, hedge fund, banque} \iqlevel{junior, confirmé}]\label{iq:iv2:vol-surface-and-vol-trading:12}
Le cours de l'action étant de 100, la volatilité implicite à un an à la monnaie est de 20~\%, et le smile perd un point pour chaque hausse de 10~\% du strike. L'action bondit à 105. Quelle est la volatilité à la monnaie, et quelle est la volatilité du strike 100, en sticky strike puis en sticky moneyness~? Vous êtes acheteur du call 100, avec un vega de 0.39 par point~: quelle est la différence sur votre profit~?
\end{question}

\begin{question}[$\star\star$ \iqroles{QT, Trader, QR, Risque} \iqfirm{teneur de marché, hedge fund, banque} \iqlevel{junior, confirmé}]\label{iq:iv2:vol-surface-and-vol-trading:13}
Un call à un an à la monnaie sur une action à 100~\$ présente un delta de Black--Scholes de 0.54 et un vega de 39.7 par unité de volatilité, et la pente du smile vaut $-0.002$ par dollar de strike. Quel delta utilisez-vous en sticky strike, en sticky moneyness et avec un modèle de volatilité locale qui déplace le smile dans l'autre sens~? Lequel est le bon~?
\end{question}

\begin{question}[$\star\star$ \iqroles{QT, Trader} \iqfirm{teneur de marché, hedge fund, banque} \iqlevel{junior, confirmé}]\label{iq:iv2:vol-surface-and-vol-trading:14}
Le smile à un mois a une pente de $-0.5$ point de volatilité pour 1~\% de strike. Une action recule de 2~\% et la volatilité implicite à la monnaie monte de 1.5 point. Quelle part de cela correspond au glissement de l'action le long du smile, et quelle part au déplacement du smile lui-même~? Que s'est-il passé pour l'option dont le strike était fixé à l'ancien niveau à la monnaie~?
\end{question}

\begin{question}[$\star\star$ \iqroles{QT, Trader, QR} \iqfirm{hedge fund, teneur de marché} \iqlevel{junior, confirmé}]\label{iq:iv2:vol-surface-and-vol-trading:15}
Vous vendez un swap de variance à un an sur un indice avec un strike de 22, avec un notionnel vega de 100\,000~\$, et sa réalisation s'établit à 18. Combien gagnez-vous~? À peu près combien de points de vega cela représente-t-il~? Pourquoi la volatilité implicite reste-t-elle en moyenne au-dessus de la volatilité réalisée, et quel est le piège de cette vente~?
\end{question}

\begin{question}[$\star\star$ \iqroles{QR, QT} \iqfirm{hedge fund, teneur de marché} \iqlevel{junior, confirmé}]\label{iq:iv2:vol-surface-and-vol-trading:16}
Vous disposez de données d'ouverture, de plus haut, de plus bas et de clôture. Comment estimez-vous la variance d'une journée à partir de l'amplitude~? D'où vient le facteur $1/(4\ln 2)$, et de combien est-il meilleur qu'un rendement close-to-close au carré~? Qu'est-ce qui le biaise sur des données réelles~?
\end{question}

\begin{question}[$\star\star$ \iqroles{QT, Trader, QR} \iqfirm{teneur de marché, hedge fund} \iqlevel{junior, confirmé}]\label{iq:iv2:vol-surface-and-vol-trading:17}
Nous sommes à la clôture du vendredi. Une action affiche une volatilité en temps ouvré de 25~\%, comptée sur 252 jours de cotation. Une option expire à la clôture du lundi, soit à un jour de cotation et trois jours calendaires. Quelle volatilité un écran affiche-t-il, en comptant 365 jours calendaires~? Pourquoi paraît-elle si faible~? Qu'est-ce que cela change pour un trader qui regarde la structure par terme~?
\end{question}

\begin{question}[$\star\star$ \iqroles{QT, Trader, QR, Analyste} \iqfirm{teneur de marché, hedge fund} \iqlevel{junior, confirmé}]\label{iq:iv2:vol-surface-and-vol-trading:18}
La publication des résultats tombe dans les cinq prochains jours de cotation. L'option à cinq jours a une volatilité implicite de 60~\%, et l'option à 25 jours de 40~\%. Aucune des deux ne comporte d'autre événement. Quelle est la volatilité de base et quel est le mouvement implicite le jour de l'événement, exprimé en écart type~? Quel est le mouvement absolu attendu~? Quelles hypothèses faites-vous~?
\end{question}

\begin{question}[$\star\star$ \iqroles{QT, Trader, Analyste} \iqfirm{teneur de marché, hedge fund} \iqlevel{junior, confirmé}]\label{iq:iv2:vol-surface-and-vol-trading:19}
Une action à 100~\$ a une volatilité de base de 25~\% et un mouvement lié aux résultats d'un écart type de 6~\% demain, avec 30 jours calendaires avant l'échéance, sur une année de 365 jours. Quelle est la volatilité implicite de l'option à la monnaie~? Vous achetez le straddle à environ 7.46~\$. Après l'événement, la volatilité implicite est de 25~\%, et il reste 29 jours. Combien vaut le straddle si l'action reste inchangée, ou si elle se situe à 103, 106 ou 110~? De quelle distance l'action doit-elle se déplacer pour que vous atteigniez le point mort~?
\end{question}

\begin{question}[$\star\star$ \iqroles{QT, QR} \iqfirm{hedge fund, teneur de marché} \iqlevel{junior, confirmé}]\label{iq:iv2:vol-surface-and-vol-trading:20}
La volatilité réalisée de l'année prochaine a une moyenne de 20~\% et un écart type de 5~\%. Quel est le strike juste du swap de variance~? Quel est le strike juste du swap de volatilité~? Pourquoi un swap de volatilité vous verse-t-il moins lors d'un grand mouvement~? Pourquoi est-il difficile à couvrir~?
\end{question}

\begin{question}[$\star\star$ \iqroles{QT, QR, Trader} \iqfirm{hedge fund, teneur de marché, banque} \iqlevel{junior, confirmé}]\label{iq:iv2:vol-surface-and-vol-trading:21}
Un indice compte 50 actions pondérées à parts égales, chacune avec une volatilité de 30~\% et une corrélation par paire de 0.3. Quelle est la volatilité de l'indice~? Vous vendez la variance de l'indice et achetez la variance de toutes les actions. Si la corrélation se réalise à 0.2 et la volatilité des actions à 30~\%, en quoi cela vous avantage-t-il~? Et si la volatilité des actions vaut 33~\%~? Comment s'appelle cette opération, et quel risque portez-vous réellement~?
\end{question}

\begin{question}[$\star\star$ \iqroles{QT, Trader, QR} \iqfirm{teneur de marché, hedge fund, banque} \iqlevel{junior}]\label{iq:iv2:vol-surface-and-vol-trading:22}
Pour une échéance d'un an sur une action à 100~\$, la volatilité à la monnaie est de 20~\% et celle du strike 90 de 25~\%. Combien coûte un put spread 100-90~? Combien coûterait-il si vous évaluiez les deux options à 20~\%~? Pourquoi le spread est-il moins cher avec le skew~? Qui paie donc quoi~?
\end{question}

\begin{question}[$\star\star$ \iqroles{QR, QD} \iqfirm{hedge fund, banque, teneur de marché} \iqlevel{junior, confirmé}]\label{iq:iv2:vol-surface-and-vol-trading:23}
Un modèle de Heston a une vitesse de retour à la moyenne de la variance de 2, une variance de long terme de 0.04 et une volatilité de la volatilité de 0.5. Satisfait-il la condition de Feller~? Que signifie-t-il s'il ne la satisfait pas, et que faites-vous quand vous le simulez~?
\end{question}

\begin{question}[$\star\star$ \iqroles{QR, QD, Risque} \iqfirm{toutes} \iqlevel{junior, confirmé}]\label{iq:iv2:vol-surface-and-vol-trading:24}
Vous devez interpoler la volatilité à la monnaie à deux mois, sachant que celle à un mois vaut 20~\% et celle à trois mois 22~\%. Comment interpolez-vous~? Quel résultat obtenez-vous~? Pourquoi ne pas interpoler linéairement en volatilité~? Comment vous assurez-vous qu'aucun arbitrage calendaire n'apparaît~?
\end{question}

\begin{question}[$\star\star$ \iqroles{QT, QR, Risque} \iqfirm{hedge fund, teneur de marché, banque} \iqlevel{junior, confirmé}]\label{iq:iv2:vol-surface-and-vol-trading:25}
Vous êtes acheteur d'un swap de variance à un an avec un notionnel vega de 100\,000~\$ et un strike de 20. Après six mois, la volatilité réalisée jusqu'ici est de 25, et la volatilité implicite pour la période restante est de 22. Quelle est sa valeur~? Pourquoi cette valorisation est-elle si simple~? Comment son vega se compare-t-il à celui du départ~?
\end{question}

\section{Variance et modèles}

\begin{question}[$\star\star\star$ \iqroles{QR, QT} \iqfirm{hedge fund, teneur de marché} \iqlevel{junior, confirmé}]\label{iq:iv2:vol-surface-and-vol-trading:26}
Pourquoi la réplication d'un swap de variance utilise-t-elle des options pondérées par l'inverse du carré du strike~? Montrez que, pour une action à 80, 100 et 120, avec une volatilité plate de 20~\% sur un an, ce strip a le même gamma en dollars en chacun de ces niveaux, et qu'un strip à poids égaux ne l'a pas. Quel est le strike juste, et quel rapport cela a-t-il avec le contrat log~?
\end{question}

\begin{question}[$\star\star\star$ \iqroles{QR} \iqfirm{hedge fund, banque} \iqlevel{confirmé}]\label{iq:iv2:vol-surface-and-vol-trading:27}
Le smile à un an d'un fournisseur de données monte fortement dans les ailes lointaines, et sa variance totale vaut $3|\ln(K/S)|$ pour les grands strikes. Montrez à partir du prix du call que c'est impossible. Quelle est la pente maximale admise~? Pour une action à 100, quelle est la volatilité implicite maximale à un strike de 400, si la borne s'applique~?
\end{question}

\begin{question}[$\star\star\star$ \iqroles{QR, QD} \iqfirm{hedge fund, banque, teneur de marché} \iqlevel{junior, confirmé}]\label{iq:iv2:vol-surface-and-vol-trading:28}
Écrivez la formule de Dupire pour la volatilité locale avec un taux nul. Vérifiez-la pour une surface plate à 20~\%. Vérifiez-la ensuite pour le mélange dans lequel les calls sont la moyenne des prix de Black-Scholes à 10~\% et à 30~\%, pour un strike à un an de 100. Quel résultat obtenez-vous~? Pourquoi est-il si différent de la volatilité implicite de 19.98~\%, et comment l'interpréteriez-vous~?
\end{question}

\begin{question}[$\star\star\star$ \iqroles{QR, QT} \iqfirm{hedge fund, teneur de marché, banque} \iqlevel{junior, confirmé}]\label{iq:iv2:vol-surface-and-vol-trading:29}
Une action à 100~\$ a une volatilité locale de $20~\% - 0.2\ln(S/100)$, avec un taux nul. Quelle est la volatilité implicite à un an pour des strikes de 80, 100 et 120~? Comparez la pente implicite en log-strike à la pente locale. Quelle est la règle, et pourquoi fonctionne-t-elle~?
\end{question}

\begin{question}[$\star\star\star$ \iqroles{QR, QT} \iqfirm{hedge fund, teneur de marché, banque} \iqlevel{confirmé}]\label{iq:iv2:vol-surface-and-vol-trading:30}
Une action a une volatilité locale de $20~\% - 0.2\ln(S/100)$, avec un taux nul. L'action valant 100, la volatilité implicite à un an pour un strike de 100 vaut 19.98~\%. L'action valant maintenant 105, elle vaut 19.01~\% pour un strike de 105 et 19.48~\% pour un strike de 100. Que fait la volatilité à la monnaie, et dans quel sens le smile s'est-il déplacé~? Comment cela se compare-t-il à ce que font les marchés en général~? Qu'est-ce que cela fait à votre couverture~?
\end{question}

\begin{question}[$\star\star\star$ \iqroles{QR, QD} \iqfirm{hedge fund, banque, teneur de marché} \iqlevel{junior, confirmé}]\label{iq:iv2:vol-surface-and-vol-trading:31}
Dans le modèle de Heston, la variance actuelle vaut 0.04 et son niveau de long terme 0.09, avec une vitesse de retour de 2. À peu près, à quoi ressemble la structure par terme à la monnaie à trois mois, un an et quatre ans, avant tout effet de la volatilité de la volatilité~? Fixez maintenant la variance et le niveau de long terme tous deux à 0.04, gardez la vitesse de 2, et prenez une volatilité de la volatilité de 0.5, sur une action à 100~\$ avec un taux nul. La volatilité implicite à un an, en pourcentage, pour des strikes de 80, 100 et 120 vaut 23.1, 18.1 et 14.6 pour une corrélation de $-0.7$, et 20.4, 18.8 et 19.9 pour une corrélation nulle. Qu'est-ce qui fixe la pente, et qu'est-ce qui fixe la courbure~? Pourquoi la volatilité à la monnaie est-elle inférieure à 20~\%~?
\end{question}

\begin{question}[$\star\star\star$ \iqroles{QR, QT} \iqfirm{banque, hedge fund} \iqlevel{junior, confirmé}]\label{iq:iv2:vol-surface-and-vol-trading:32}
Dans le modèle SABR lognormal, avec une volatilité initiale de 20~\%, une volatilité de la volatilité de 0.4, une corrélation de $-0.3$ et un an, quelle est la volatilité implicite à la monnaie donnée par la formule de Hagan~? Que donne-t-elle pour des strikes de 80 et 120~? Comment la vérifieriez-vous~? Quel est l'avantage du modèle SABR sur la volatilité locale, et quelle est sa faiblesse bien connue~?
\end{question}

\begin{question}[$\star\star\star$ \iqroles{QR, QT} \iqfirm{hedge fund, teneur de marché, banque} \iqlevel{junior, confirmé}]\label{iq:iv2:vol-surface-and-vol-trading:33}
Une action à 100~\$, avec un taux nul, a une volatilité de diffusion de 15~\%, et elle peut chuter brutalement d'un saut à la baisse de 20~\%, à raison de 0.1 fois par an. Écrivez son prix de call comme un mélange de Poisson de prix de Black-Scholes, avec la bonne dérive. Quelles sont les volatilités implicites pour des strikes de 90, 100 et 110, à un mois et à un an~? Pourquoi le skew à court terme est-il si raide~? Qu'est-ce qu'un modèle de diffusion pur ne peut pas faire ici~?
\end{question}

\begin{question}[$\star\star\star$ \iqroles{QR, QT} \iqfirm{hedge fund, teneur de marché} \iqlevel{confirmé}]\label{iq:iv2:vol-surface-and-vol-trading:34}
Vous avez vendu un swap de variance et vous l'avez couvert avec le strip de contrats log. L'action saute de 20~\% vers le bas en un seul pas, ou de 20~\% vers le haut. Que verse le rendement log au carré du swap dans chaque cas, et que rapporte votre couverture par le strip~? Qui perd, et de combien par unité de variance~? Qu'est-ce que cela implique pour le strike~?
\end{question}

\begin{question}[$\star\star\star$ \iqroles{QR, QT} \iqfirm{hedge fund, teneur de marché} \iqlevel{junior, confirmé}]\label{iq:iv2:vol-surface-and-vol-trading:35}
Un smile à un an est donné par $\sigma(K) = 20~\% - 0.1\ln(K/100)$, avec un plancher de 3~\%, pour une action à 100~\$ avec un taux nul. À l'aide du strip d'options hors de la monnaie, quel est le strike juste du swap de variance~? La règle d'un desk est $\sigma_{\text{ATM}}^2(1 + 3T\,b^2)$, avec une pente $b$ par unité de log-strike. Que donne cette règle~? Pourquoi le strike est-il supérieur à la volatilité à la monnaie~?
\end{question}

\begin{question}[$\star\star\star$ \iqroles{QT, QR} \iqfirm{hedge fund, teneur de marché, banque} \iqlevel{junior, confirmé}]\label{iq:iv2:vol-surface-and-vol-trading:36}
Les actions ont une volatilité de 30~\%. Un indice composé d'un très grand nombre d'entre elles vaut 100, avec des taux nuls. Sa corrélation sur l'année vaut 0.2 avec une probabilité de 4/5 et 0.8 avec une probabilité de 1/5. Quelle est la volatilité de l'indice dans chaque état~? Quel est son smile à un an pour des strikes de 80, 100 et 120, si les deux états ont le même forward~? Cela donne-t-il un skew~? Quel ingrédient du modèle faut-il ajouter pour en obtenir un, et pourquoi les skews d'indice sont-ils plus raides que ceux sur une seule action~?
\end{question}

\begin{question}[$\star\star\star$ \iqroles{QT, Trader, QR} \iqfirm{hedge fund, teneur de marché} \iqlevel{junior, confirmé}]\label{iq:iv2:vol-surface-and-vol-trading:37}
Vous vendez un straddle à la monnaie à un mois sur une action à 100~\$, à une volatilité de 20~\%, alors que la volatilité de l'action vaut en réalité 16~\%. Vous couvrez chaque jour, soit 21 fois. Quelles sont la moyenne et l'écart type de votre profit~? À quelle fréquence perdez-vous, et quel est le 1\up{er}~centile~? Que dit cela de la taille des positions à prendre, lorsque les distributions réelles comportent des sauts~?
\end{question}

\section{Trois à construire}

\begin{question}[$\star\star\star\star$ \iqroles{QR, QD} \iqfirm{hedge fund, teneur de marché} \iqlevel{confirmé}]\label{iq:iv2:vol-surface-and-vol-trading:38}
Dérivez le strike juste du swap de variance à partir du contrat log, et écrivez-le comme un strip d'options hors de la monnaie. Calculez-le pour le smile $\sigma(K) = 20~\% - 0.1\ln(K/100)$, avec un plancher de 3~\%, sur une action à 100~\$ avec un taux nul sur un an, en utilisant des strikes de 50 à 200 par pas de 1~\$, 5~\$ et 10~\$, puis de 70 à 140 par pas de 5~\$. Comparez-le à l'intégrale exacte. À quoi tient chaque erreur~? Laquelle compte le plus en situation de stress~?
\end{question}

\begin{question}[$\star\star\star\star$ \iqroles{QR, QD} \iqfirm{hedge fund, banque} \iqlevel{confirmé}]\label{iq:iv2:vol-surface-and-vol-trading:39}
Montrez que, si la volatilité est aléatoire et indépendante du bruit propre de l'action, le prix d'un call est la moyenne des prix de Black-Scholes sur la distribution de la variance réalisée. Testez-le sur un modèle de Heston pour une action à 100~\$ sur un an avec un taux nul, $v_0 = \theta = 0.04$, $\kappa = 2$, une volatilité de la volatilité de 0.5 et une corrélation nulle, en le comparant à une simulation classique du même modèle pour des strikes de 80, 100 et 120. De combien diminue l'erreur type~? Que devient l'argument en présence de corrélation~?
\end{question}

\begin{question}[$\star\star\star\star$ \iqroles{QR, QD} \iqfirm{hedge fund, banque} \iqlevel{confirmé}]\label{iq:iv2:vol-surface-and-vol-trading:40}
Une action à 100~\$ avec un taux nul a, pour chaque strike et chaque échéance, un prix de call égal à la moyenne à poids égaux des prix de Black-Scholes à des volatilités de 10~\% et de 30~\%. Trouvez la surface de volatilité locale sous forme fermée, comme une moyenne des deux variances, les poids étant donnés par la densité de chaque état. Recalculez les calls à un an pour des strikes de 70, 100 et 130 dans ce modèle de volatilité locale. Une simulation d'Euler quotidienne simple retrouve-t-elle ces prix~? Qu'est-ce qui se passe mal au voisinage de l'instant zéro, et comment le résolvez-vous correctement~?
\end{question}
